% META: {"id":"1SPE-SUITES-FR-R3","chapitre":"1SPE-SUITES","type_objet":"remediation","status":"approved","prerequis_testes":["R3"],"origin":"generated","status_history":["generated","audited","accepted","approved"],"release_acceptance":"RELEASE_OWNER_BATCH_ACCEPTANCE","acceptance_closure_digest":"sha256:9b3ccf9a81c5520fbb7e03b7b2d3e7bf2057a02834b6d908bf3be5f49f3b553f","acceptance_receipt_digest":"sha256:4fad86f0282e0fa4e0419cca1ad6379ae7af5a280c04a4a90880f90a1c044242"}

%---------------------------------------------------------------
% FICHE DE REMISE À NIVEAU — R3 : Fonctions
% Image, antécédent, sens de variation
%---------------------------------------------------------------

\subsection*{Rappel — Fonctions : image, antécédent, variations}

Soit $f$ une fonction définie sur un intervalle $I$.

\textbf{Image :} $f(a)$ est l'\emph{image} de $a$ par $f$ (on calcule en substituant $x = a$).

\textbf{Antécédent :} $a$ est un \emph{antécédent} de $b$ par $f$ si $f(a) = b$ (on résout l'équation $f(x) = b$).

\textbf{Sens de variation :}
\begin{itemize}
  \item $f$ est \textbf{croissante} sur $I$ si, pour tous $x_1 < x_2$ dans $I$, $f(x_1) \leq f(x_2)$.
  \item $f$ est \textbf{décroissante} sur $I$ si, pour tous $x_1 < x_2$ dans $I$, $f(x_1) \geq f(x_2)$.
\end{itemize}

\textbf{Lien avec les suites :} une suite $(u_n)$ définie par $u_n = f(n)$ hérite des variations de $f$ sur $\mathbb{N}$.

%---------------------------------------------------------------

\begin{exercice}{1SPE-SUITES-FR-R3-EX1}{1}{5}
On considère la fonction $f$ définie sur $\mathbb{R}$ par $f(x) = 2x^2 - 3x + 1$.
\begin{enumerate}
  \item Calculer $f(0)$, $f(1)$, $f(2)$ et $f(3)$.
  \item Vérifier que $1$ est un antécédent de $0$ par $f$.
  \item Trouver tous les antécédents de $0$ par $f$.
  \item La suite $(u_n)$ définie par $u_n = f(n)$ pour $n \geq 2$ est-elle croissante ?
\end{enumerate}
\end{exercice}

% BEGIN-VERIFY
% from sympy import symbols, solve, factor, Rational
% x = symbols('x')
% f = 2*x**2 - 3*x + 1
% assert f.subs(x, 0) == 1
% assert f.subs(x, 1) == 0
% assert f.subs(x, 2) == 3
% assert f.subs(x, 3) == 10
% sols = solve(f, x)
% assert set(sols) == {1, Rational(1, 2)}
% # Pour n >= 2, f est croissante (sommet en x=3/4 < 2)
% assert f.subs(x, 2) < f.subs(x, 3)
% assert f.subs(x, 3) < f.subs(x, 4)
% END-VERIFY

\begin{corrige}{1SPE-SUITES-FR-R3-EX1}

\commentaireMarge{On substitue les valeurs de $n$ dans la formule de $f$.}

\textbf{Question 1.}
\[
f(0) = 2 \times 0 - 3 \times 0 + 1 = 1.
\]
\[
f(1) = 2 \times 1 - 3 \times 1 + 1 = 2 - 3 + 1 = 0.
\]
\[
f(2) = 2 \times 4 - 3 \times 2 + 1 = 8 - 6 + 1 = 3.
\]
\[
f(3) = 2 \times 9 - 3 \times 3 + 1 = 18 - 9 + 1 = 10.
\]

\commentaireMarge{Comparer l'image à la cible.}

\textbf{Question 2.} On calcule $f(1) = 0$ (voir question 1). Donc $f(1) = 0$ : $1$ est bien un antécédent de $0$.

\textbf{Question 3.} On résout $2x^2 - 3x + 1 = 0$. On factorise :
\[
2x^2 - 3x + 1 = (2x - 1)(x - 1) = 0.
\]
Donc $x = \dfrac{1}{2}$ ou $x = 1$. Les antécédents de $0$ sont $\dfrac{1}{2}$ et $1$.

\commentaireMarge{$f$ croissante à droite de $x_s = 3/4$.}

\textbf{Question 4.} Le sommet de la parabole est en $x_s = \dfrac{3}{2 \times 2} = \dfrac{3}{4}$. Pour $n \geq 2 > \dfrac{3}{4}$, la fonction $f$ est croissante, donc la suite $(u_n)_{n \geq 2}$ est croissante.

\end{corrige}

%---------------------------------------------------------------

\begin{exercice}{1SPE-SUITES-FR-R3-EX2}{1}{5}
On considère la fonction $g$ définie sur $\mathbb{R}^+$ par $g(x) = \dfrac{x + 4}{x + 1}$.
\begin{enumerate}
  \item Calculer $g(0)$, $g(1)$, $g(3)$ et $g(9)$.
  \item Montrer que, pour tout $x \geq 0$, $g(x) - 1 = \dfrac{3}{x+1}$.
  \item En déduire le sens de variation de $g$ sur $[0\,;+\infty[$.
\end{enumerate}
\end{exercice}

% BEGIN-VERIFY
% from sympy import symbols, simplify, Rational
% x = symbols('x', nonneg=True)
% g = (x + 4) / (x + 1)
% assert g.subs(x, 0) == 4
% assert g.subs(x, 1) == Rational(5, 2)
% assert g.subs(x, 3) == Rational(7, 4)
% assert g.subs(x, 9) == Rational(13, 10)
% diff = simplify(g - 1 - 3/(x+1))
% assert diff == 0
% END-VERIFY

\begin{corrige}{1SPE-SUITES-FR-R3-EX2}

\textbf{Question 1.}
\[
g(0) = \frac{0+4}{0+1} = 4. \quad g(1) = \frac{1+4}{1+1} = \frac{5}{2} = 2{,}5.
\]
\[
g(3) = \frac{3+4}{3+1} = \frac{7}{4} = 1{,}75. \quad g(9) = \frac{9+4}{9+1} = \frac{13}{10} = 1{,}3.
\]

\commentaireMarge{On met au même dénominateur en écrivant $1 = \frac{x+1}{x+1}$.}

\textbf{Question 2.}
\[
g(x) - 1 = \frac{x+4}{x+1} - 1 = \frac{x+4}{x+1} - \frac{x+1}{x+1} = \frac{(x+4)-(x+1)}{x+1} = \frac{3}{x+1}.
\]

\commentaireMarge{Sur $[0\,;+\infty[$, quand $x$ augmente, $x+1$ augmente, donc $\frac{3}{x+1}$ diminue, donc $g(x) = 1 + \frac{3}{x+1}$ diminue.}

\textbf{Question 3.} Pour $x \geq 0$, $x+1 \geq 1 > 0$ et est croissant, donc $\dfrac{3}{x+1}$ est décroissant. Ainsi $g(x) = 1 + \dfrac{3}{x+1}$ est décroissante sur $[0\,;+\infty[$.

\end{corrige}

%---------------------------------------------------------------

\begin{exercice}{1SPE-SUITES-FR-R3-EX3}{1}{5}
On considère la suite $(v_n)$ définie pour tout $n \in \mathbb{N}$ par $v_n = 3n - 7$.
\begin{enumerate}
  \item Calculer $v_0$, $v_1$, $v_2$, $v_5$.
  \item À quel rang $n$ a-t-on $v_n = 0$ ?
  \item À partir de quel rang $n$ a-t-on $v_n > 0$ ?
  \item La suite $(v_n)$ est-elle croissante ? Justifier brièvement.
\end{enumerate}
\end{exercice}

% BEGIN-VERIFY
% from sympy import symbols, solve, Rational
% n = symbols('n')
% v = 3*n - 7
% assert v.subs(n, 0) == -7
% assert v.subs(n, 1) == -4
% assert v.subs(n, 2) == -1
% assert v.subs(n, 5) == 8
% sol = solve(v, n)
% assert sol == [Rational(7, 3)]
% # v_n > 0 iff n > 7/3 iff n >= 3
% assert v.subs(n, 3) > 0
% assert v.subs(n, 2) < 0
% END-VERIFY

\begin{corrige}{1SPE-SUITES-FR-R3-EX3}

\textbf{Question 1.}
\[
v_0 = 3 \times 0 - 7 = -7. \quad v_1 = 3 \times 1 - 7 = -4. \quad v_2 = 3 \times 2 - 7 = -1. \quad v_5 = 3 \times 5 - 7 = 8.
\]

\textbf{Question 2.} On résout $v_n = 0$ :
\[
3n - 7 = 0 \implies n = \frac{7}{3} \approx 2{,}33.
\]
Mais $n$ doit être un entier naturel : $\dfrac{7}{3}$ n'est pas entier, donc il n'existe pas de rang $n \in \mathbb{N}$ tel que $v_n = 0$.

\commentaireMarge{$n > 7/3$ et $n$ entier signifie $n \geq 3$.}

\textbf{Question 3.} $v_n > 0 \iff 3n - 7 > 0 \iff n > \dfrac{7}{3}$. Le plus petit entier vérifiant cette condition est $n = 3$. Pour tout $n \geq 3$, $v_n > 0$.

\textbf{Question 4.} $v_{n+1} - v_n = [3(n+1) - 7] - [3n - 7] = 3 > 0$. La suite est croissante.

\end{corrige}

%---------------------------------------------------------------

\begin{exercice}{1SPE-SUITES-FR-R3-EX4}{1}{6}
On considère la fonction affine $h$ définie sur $\mathbb{R}$ par
$h(x)=-2x+5$.
\begin{enumerate}
  \item Calculer $h(-1)$, $h(0)$, $h(1)$ et $h(2)$.
  \item Résoudre $h(x)=1$.
  \item Déterminer le sens de variation de $h$ sur $\mathbb{R}$.
  \item On pose $w_n=h(n)$ pour tout $n\in\mathbb{N}$. Calculer $w_0$,
  $w_1$, $w_2$, puis justifier le sens de variation de $(w_n)$.
\end{enumerate}
\end{exercice}

% BEGIN-VERIFY
% from sympy import symbols, solve
% x = symbols('x')
% h = -2*x + 5
% assert h.subs(x, -1) == 7
% assert h.subs(x, 0) == 5
% assert h.subs(x, 1) == 3
% assert h.subs(x, 2) == 1
% assert solve(h - 1, x) == [2]
% w = lambda n: -2*n + 5
% assert [w(n) for n in range(3)] == [5, 3, 1]
% assert all(w(n+1) - w(n) == -2 for n in range(5))
% END-VERIFY

\begin{corrige}{1SPE-SUITES-FR-R3-EX4}

\textbf{Question 1.}
\[
h(-1)=7. \quad h(0)=5. \quad h(1)=3. \quad h(2)=1.
\]

\textbf{Question 2.}
\[
h(x)=1 \iff -2x+5=1 \iff -2x=-4 \iff x=2.
\]

\commentaireMarge{Le coefficient directeur de $h$ vaut $-2<0$.}

\textbf{Question 3.} La fonction affine $h$ a pour coefficient directeur
$-2$, donc elle est strictement décroissante sur $\mathbb{R}$.

\textbf{Question 4.} On a $w_0=5$, $w_1=3$ et $w_2=1$. De plus,
pour tout $n\in\mathbb{N}$,
\[
w_{n+1}-w_n=[-2(n+1)+5]-[-2n+5]=-2<0.
\]
La suite $(w_n)$ est donc strictement décroissante.

\end{corrige}
